\chapter*{Tóm tắt}
\label{summary}

Quy trình xét duyệt bài báo tại hội nghị khoa học gồm nhiều công đoạn liên tiếp, từ cấu hình hội nghị, tiếp nhận bài nộp, phân công và theo dõi phản biện đến tổng hợp đánh giá và ra quyết định. Khi số bài nộp tăng, khối lượng thao tác và thông tin cần xử lý ở từng vai trò cũng tăng theo. Các nền tảng phổ biến hiện nay đã hỗ trợ tốt các nghiệp vụ tổ chức cơ bản. Trong phạm vi khảo sát của đề tài, khoảng trống chủ yếu nằm ở việc bổ sung các cơ chế hỗ trợ có kiểm soát cho một số công đoạn nặng về nhập liệu, đọc hiểu và tổng hợp thông tin, đồng thời vẫn giữ quyết định học thuật cuối cùng cho con người.

Đề tài xây dựng ConferenceSpace như một nền tảng web phục vụ vòng đời xét duyệt bài báo cho ba vai trò chính: tác giả, phản biện viên và chủ tọa. Hệ thống bao gồm các công đoạn nghiệp vụ cốt lõi như cấu hình hội nghị và theo dõi tiến độ, nộp và quản lý bản thảo, phân công phản biện, thu thập đánh giá, giai đoạn phản hồi của tác giả (rebuttal) và hỗ trợ ra quyết định. Bên cạnh đó, hệ thống có các cơ chế thuật toán có thể kiểm chứng để đề xuất phản biện viên và hỗ trợ phát hiện xung đột lợi ích, kết quả được trình bày dưới dạng gợi ý để chủ tọa xem xét. Ở một số công đoạn, hệ thống bổ sung hỗ trợ AI có kiểm soát nhằm giảm thao tác thủ công hoặc hỗ trợ đọc bằng cách tóm tắt và tổng hợp bằng chứng; người dùng có toàn quyền đưa ra quyết định, hệ thống không tự đưa ra và xác nhận quyết định thông qua các kết quả của AI.

Phần đánh giá xem xét nhiều lớp bằng chứng trong điều kiện thử nghiệm: độ ổn định của quy trình nghiệp vụ, chất lượng đề xuất phân công, và mức độ sử dụng được của các luồng hỗ trợ AI. Chức năng đề xuất phản biện đưa chuyên gia phù hợp vào 10 vị trí đầu trong khoảng 65\% trường hợp thử và đáp ứng khoảng 65,9\% nhu cầu phân công trên tập thử. Trên tập 1.127 bài báo thật, hỗ trợ điền thông tin bài nộp đạt độ chính xác khoảng 98,20\% với tiêu đề và 92,77\% với từ khóa; phân tích bài báo đạt khoảng 96,22\% độ trung thực trích dẫn; tổng hợp bằng chứng hỗ trợ chủ tọa đạt khoảng 87,34\% độ trung thực theo tiêu chí thử nghiệm. Đồng thời, chỉ khoảng 47\% cảnh báo về nội dung phản biện có thể được chứng minh đầy đủ, cho thấy kết quả do AI tạo ra còn giới hạn và cần người dùng kiểm tra. Kết quả cho thấy ConferenceSpace có thể vận hành như một hệ thống xét duyệt quy mô thử nghiệm, trong đó AI chỉ là lớp hỗ trợ bổ sung ở một số công đoạn, còn quyền đánh giá và quyết định học thuật vẫn thuộc về con người.